% Typography and page geometry
\usepackage[margin=30truemm]{geometry}
\usepackage{stix2}
\usepackage{microtype}
\usepackage{amsmath,amsthm,mathtools,amscd}
\usepackage{aliascnt}
\usepackage{enumitem}
\usepackage{xcolor}
\usepackage{tikz}
\usepackage{tikz-cd}
\usetikzlibrary{positioning,arrows.meta,decorations.markings,calc}

% Bibliography
\usepackage[backend=biber,style=alphabetic,sorting=nyt,maxnames=20]{biblatex}
\addbibresource{my.bib}

% Hyperlinks and references
\usepackage{xurl}
\usepackage{hyperref}
\usepackage{bookmark}
\usepackage{cleveref}

\definecolor{FushimiBlue}{HTML}{175A8C}
\definecolor{FushimiRed}{HTML}{8B2F3F}
\definecolor{FushimiViolet}{HTML}{66507A}
\definecolor{FushimiOrange}{HTML}{A25B23}

\hypersetup{
	colorlinks=true,
	linkcolor=FushimiRed,
	citecolor=FushimiOrange,
	urlcolor=FushimiOrange,
	filecolor=FushimiBlue,
	pdfauthor={Riku Fushimi},
	pdftitle={Cluster Morita theorem for negative cluster categories}
}

% Paragraphs and lists
\allowdisplaybreaks
\setlength{\parindent}{1.45em}
\setlength{\parskip}{0pt}
\setlength{\emergencystretch}{2em}
\setlist[enumerate]{leftmargin=2.35em,itemsep=0.12em,topsep=0.32em,parsep=0pt}
\setlist[itemize]{leftmargin=2.15em,itemsep=0.12em,topsep=0.32em,parsep=0pt}

% Theorem environments
\newtheoremstyle{fushimi}
	{0.65\baselineskip}
	{0.55\baselineskip}
	{\normalfont}
	{}
	{\bfseries}
	{.}
	{0.5em}
	{}
\theoremstyle{fushimi}
\newtheorem{mainthm}{Theorem}
\renewcommand{\themainthm}{\Alph{mainthm}}
\newtheorem{thm}{Theorem}[section]

\newaliascnt{prop}{thm}
\newtheorem{prop}[prop]{Proposition}
\aliascntresetthe{prop}
\newaliascnt{lem}{thm}
\newtheorem{lem}[lem]{Lemma}
\aliascntresetthe{lem}
\newaliascnt{cor}{thm}
\newtheorem{cor}[cor]{Corollary}
\aliascntresetthe{cor}
\newaliascnt{dfn}{thm}
\newtheorem{dfn}[dfn]{Definition}
\aliascntresetthe{dfn}
\newaliascnt{rmk}{thm}
\newtheorem{rmk}[rmk]{Remark}
\aliascntresetthe{rmk}
\newaliascnt{exa}{thm}
\newtheorem{exa}[exa]{Example}
\aliascntresetthe{exa}

\numberwithin{equation}{section}

\crefname{mainthm}{Theorem}{Theorems}
\Crefname{mainthm}{Theorem}{Theorems}
\crefname{thm}{Theorem}{Theorems}
\crefname{prop}{Proposition}{Propositions}
\crefname{lem}{Lemma}{Lemmas}
\crefname{cor}{Corollary}{Corollaries}
\crefname{dfn}{Definition}{Definitions}
\crefname{rmk}{Remark}{Remarks}
\crefname{exa}{Example}{Examples}
\crefname{section}{Section}{Sections}
\crefname{subsection}{Subsection}{Subsections}

\newcommand{\defterm}[1]{\textcolor{blue}{\textit{#1}}}

% Operators
\DeclareMathOperator{\Hom}{Hom}
\DeclareMathOperator{\End}{End}
\DeclareMathOperator{\Ext}{Ext}
\DeclareMathOperator{\id}{id}
\DeclareMathOperator{\Filt}{Filt}
\DeclareMathOperator{\add}{add}
\DeclareMathOperator{\thick}{thick}
\DeclareMathOperator{\per}{per}
\DeclareMathOperator{\pvd}{pvd}
\DeclareMathOperator{\cosg}{cosg}
\DeclareMathOperator{\sg}{sg}
\DeclareMathOperator{\CM}{CM}
\DeclareMathOperator{\modu}{mod}
\DeclareMathOperator{\REnd}{\mathbf{R}End}
\DeclareMathOperator{\RHom}{\mathbf{R}Hom}
\DeclareMathOperator{\HHom}{HH}
\DeclareMathOperator{\HChom}{HC}
\DeclareMathOperator{\HN}{HN}

% Frequently used notation
\newcommand{\A}{\mathcal A}
\newcommand{\C}{\mathcal C}
\newcommand{\D}{\mathcal D}
\newcommand{\Hcal}{\mathcal H}
\newcommand{\Lcal}{\mathcal L}
\newcommand{\F}{\mathcal F}
\newcommand{\DD}{\mathbb D}
\newcommand{\Serre}{\mathbb S}
\newcommand{\Xcal}{\mathcal X}
\newcommand{\Ycal}{\mathcal Y}
\newcommand{\Scal}{\mathcal S}
\newcommand{\Pdg}{\mathcal P}
\newcommand{\Tdg}{\mathcal T}
\newcommand{\Qdg}{\mathcal Q}
\newcommand{\Ecal}{\mathcal E}
\newcommand{\Dk}{\mathbf D}
\newcommand{\op}{\mathrm{op}}
\newcommand{\dg}{\mathrm{dg}}
\newcommand{\xto}[1]{\xrightarrow{#1}}
\newcommand{\mono}{\hookrightarrow}
